\chapter{Experimental Results}
\label{ch:results}

This chapter reports the results in two parts. Section~\ref{sec:res-performance}
presents the predictive performance of the three models and
Section~\ref{sec:res-fi} presents the feature-importance rankings, whose systematic
divergence is the central finding. Everything here describes what the models rely on
to predict, read as associative rather than causal, with the interpretation left to
Chapter~\ref{ch:discussion}.

\section{Predictive Performance}
\label{sec:res-performance}

Table~\ref{tab:performance} reports each model at its best hyperparameter
configuration, averaged over the $30$ cross-validation fits with standard deviations.
Shown through the primary metric $F_2$, the two tree ensembles clearly outperform the
linear baseline. Gradient Boosting reaches the highest $F_2$ at $0.550$, the Random
Forest follows closely at $0.546$ and the Logistic Regression trails at $0.482$.
This matches the model logic of Section~\ref{sec:models}. The ensembles can
split on combinations of ratios that a linear model cannot capture and the results
indicate, that the separation between healthy and bankrupt firms is not purely linear. They
gain that edge through recall, both catching roughly $74\%$ of bankruptcies against
the Logistic Regression's $55\%$, which is exactly what the $F_2$-score rewards.
Precision is modest for all three, but the $F_2$ objective
deliberately accepts more false alarms in exchange for catching more
bankruptcies. The ROC-AUC values, all above $0.92$, look strong at first, yet 
the symmetric metric flatters the models under a $96.8/3.2$ split,
which is exactly why the selection rests on $F_2$ instead.

\begin{table}[htbp]
  \centering
  \caption{Cross-validated performance at the best configuration}
  \label{tab:performance}
  \begin{tabular}{lcccc}
    \toprule
    \textbf{Model} & \textbf{$F_2$} & \textbf{Recall} & \textbf{Precision}
      & \textbf{ROC-AUC} \\
    \midrule
    Logistic Regression & $0.482 \pm 0.117$ & $0.554 \pm 0.142$ & $0.322 \pm 0.078$
      & $0.920 \pm 0.033$ \\
    Random Forest       & $0.546 \pm 0.062$ & $0.739 \pm 0.097$ & $0.269 \pm 0.036$
      & $0.935 \pm 0.024$ \\
    Gradient Boosting   & $\mathbf{0.550 \pm 0.063}$ & $0.745 \pm 0.094$ & $0.271 \pm 0.036$
      & $0.934 \pm 0.023$ \\
    \bottomrule
  \end{tabular}
  \floatfoot{Source: own. Primary metric $F_2$; secondary metrics reported but not
    used for selection.}
\end{table}

The selected hyperparameters are listed in Table~\ref{tab:selected}. For the logistic
regression, as noted in Section~\ref{sec:lr}, the model uses Lasso rather than Ridge,
with $C = 1.0$, equivalently $\lambda = 1$. Concerning the two tree columns, the same
parameters can carry very different meanings. The number of trees,
\texttt{n\_estimators}, for example: for the random forest more trees never hurt,
they only average away more variance until performance plateaus, so the choice is
bounded by computation rather than by accuracy, whereas for gradient boosting each
added tree is another correction step that lowers bias but eventually starts to overfit, 
which makes it genuinely bias--variance relevant. \texttt{max\_depth}, as already
shown, favours deep trees for the Random Forest ($10$), which it
can afford because Bagging and feature subsampling decorrelate them, whereas Gradient
Boosting favours flat trees ($2$), each a weak learner whose bias is corrected
only step by step, so the two depths are not comparable. Two remaining choices are
worth noting. The random forest's \texttt{max\_features} settled at $0.2$, below the
usual \texttt{sqrt} or \texttt{log2} heuristics, so each split sees only a fifth of
the features, which is exactly the stronger decorrelation that $95$ partly
redundant ratios call for. Gradient Boosting settled on the smaller learning rate of
$0.05$ to reduce bias gradually over many small steps rather than a few large ones. 

\begin{table}[htbp]
  \centering
  \caption{Selected hyperparameters (highest fold-averaged $F_2$).}
  \label{tab:selected}
  \begin{tabular}{l|l|l}
    \toprule
    \textbf{Logistic Regression} & \textbf{Random Forest} & \textbf{Gradient Boosting} \\
    \midrule
    penalty $=\ell_1$ & \texttt{n\_estimators} $=300$    & \texttt{n\_estimators} $=200$    \\
    $C = 1.0$         & \texttt{max\_depth} $=10$        & \texttt{max\_depth} $=2$         \\
                      & \texttt{max\_features} $=0.2$    & \texttt{learning\_rate} $=0.05$  \\
                      & \texttt{min\_samples\_leaf} $=3$ & \texttt{min\_samples\_leaf} $=5$ \\
                      &                                  & \texttt{subsample} $=0.8$        \\
    \bottomrule
  \end{tabular}
  \floatfoot{Source: own.}
\end{table}

\section{Feature Importance Results}
\label{sec:res-fi}

The importance results are where the analysis becomes most informative. Applying the
three methods of Chapter~\ref{ch:featureimportance} to the three models does not
produce one ranking but a systematic divergence. Table~\ref{tab:divergence} shows this
by giving, for each model, its single most important feature under each method.
The pattern is read most clearly one model at a time, starting from the simple 
linear model to the two more complex ensembles, where they begin to disagree.

\begin{table}[htbp]
  \centering
  \caption{Best ranked feature under each importance method, by model. Values are the
    respective importance scores: signed coefficient for native LR; Gini impurity importance
    for native RF/GB; mean $F_2$ drop for permutation; mean $|\text{SHAP}|$ for SHAP.}
  \label{tab:divergence}
  % Scores sit on a second line below the feature name: with the values inline the
  % table runs 100pt past \textwidth, and no font reduction closes that gap.
  \begin{tabular}{llll}
    \toprule
    \textbf{Model} & \textbf{Native} & \textbf{Permutation} & \textbf{SHAP} \\
    \midrule
    Logistic Regression & Debt ratio \% & Debt ratio \%
      & Debt ratio \% \\
                        & ($+14.49$)    & ($0.248$)
      & ($0.625$) \\
    \addlinespace
    Random Forest       & Persistent EPS & Interest Expense Ratio
      & Persistent EPS \\
                        & ($0.092$)      & ($0.020$)
      & ($0.026$) \\
    \addlinespace
    Gradient Boosting   & Persistent EPS & Borrowing dependency
      & Borrowing dependency \\
                        & ($0.247$)      & ($0.024$)
      & ($0.356$) \\
    \bottomrule
  \end{tabular}
  \floatfoot{Source: own.}
\end{table}

The logistic regression is the control case with all three methods ranking the
\emph{Debt ratio} first. Its native coefficient is the largest in magnitude and
positive at $+14.49$, it takes the top permutation drop at $0.248$ and it leads the
SHAP ranking at $0.625$. This agreement is not a coincidence. Because the $\ell_1$
penalty had already thinned the redundant twins in each correlated group, little
near-identical partner information is left for the methods to disagree over, so they
largely read the same signal. It is the clearest evidence that the divergence in the
ensembles below comes from the redundancy and not from the methods being unreliable in
general. The direction also makes economic sense. The positive sign means that higher
leverage raises the predicted probability of bankruptcy. Figure~\ref{fig:perm-lr}
shows the permutation ranking, where the \emph{Debt ratio} stands clearly above the
rest. That a few ROA variants still appear indicates that the penalty thinned the redundant
group rather than removing it entirely, with the gap after ROA(B) marking it as the
feature the model leans on.

\begin{figure}[htbp]
  \centering
  \includegraphics[width=0.82\textwidth]{permutation_LogisticRegression.png}
  \caption{Permutation importance for the logistic regression, measured as the drop
    in $F_2$ when a feature is shuffled.}
  \label{fig:perm-lr}
  \floatfoot{Source: own.}
\end{figure}

The Random Forest reads differently across the three methods. Native Gini importance
and SHAP both rank \emph{Persistent EPS} first, a measure of
how steadily a firm has earned over the past four quarters, with a Gini importance of
$0.092$ and a mean SHAP value of $0.026$. Beyond that top feature SHAP spreads its
weight broadly, with no single feature standing out and most of the total sitting in
the long tail of smaller features (Figure~\ref{fig:shap-rf}).

\begin{figure}[htbp]
  \centering
  \includegraphics[width=0.82\textwidth]{shap_RandomForest_bar.png}
  \caption{Mean absolute SHAP values for the random forest.}
  \label{fig:shap-rf}
  \floatfoot{Source: own.}
\end{figure}

Permutation importance is led instead by leverage and interest measures, with the
Interest Expense Ratio first at $0.020$ and the rest of its values small
(Figure~\ref{fig:perm-rf}). Persistent EPS drops out of the ranking entirely. Its
closest earnings relatives, such as \emph{Net profit before tax/Paid-in capital} and
\emph{Per Share Net profit before tax}, are correlated with it above $0.95$, so when
it alone is shuffled they still carry almost the same information and the score barely
moves. This is the logic described in Section~\ref{sec:fi-permutation}, where
permutation hides a feature that has close substitutes and here it seems to hide the very
feature that native importance and SHAP both rank first.

\begin{figure}[htbp]
  \centering
  \includegraphics[width=0.82\textwidth]{permutation_RandomForest.png}
  \caption{Permutation importance for the random forest, drop in $F_2$ when a feature
    is shuffled.}
  \label{fig:perm-rf}
  \floatfoot{Source: own.}
\end{figure}

Gradient Boosting is the best-performing model of the three and it produces the
clearest single ranking. Under SHAP the top feature is \emph{Borrowing dependency} at
$0.356$, a direct measure of how far a firm depends on debt, which had already risen
to second place under the random forest and now leads, ahead of Persistent EPS at
$0.292$. Native Gini importance still ranks Persistent EPS first, so SHAP overturns
the impurity ordering here.

